\section{Spiders: what is proved, what is stronger, and what remains}
\label{sec:spider-specialization}

Let $G_{B,N}$ have a hub $o$ and $B$ internally disjoint arms of length $N$.
For a prefix depth $1\leq L<N$, let $U_L$ contain the hub and the first $L$
vertices of every arm; degrees are the ambient degrees of $G_{B,N}$.  This is
the first nontrivial family after a star: the quotient dimension grows with
$L$, the support can propagate along long arms, but the graph remains
bipartite and has a scalar articulation at the hub.

\subsection{Exact prefix spectrum and the SOR consequence}

The proof-owning \texttt{volume\_gated\_acceleration} note establishes
\begin{equation}
 \sigma_{U_L}=\cos\!\left(\frac{\pi}{2(L+1)}\right)
 \label{eq:spider-prefix-cross-singular}
\end{equation}
for the largest singular value in \eqref{eq:bipartite-cross-block}.  Combining
that proved spectrum with \Cref{thm:fixed-bipartite-sor} gives the following
new SOR specialization.

\begin{corollary}[Optimal SOR on an exposed spider prefix]
\label{cor:spider-prefix-sor}
On $U_L$, put
\begin{equation}
 t_L:=\sqrt{
  1-c_\alpha^2
  \cos^2\!\left(\frac{\pi}{2(L+1)}\right)
 }.
 \label{eq:spider-prefix-t}
\end{equation}
Then the optimal red--black SOR factor is
$\zeta_L=(1-t_L)/(1+t_L)$, and
\begin{equation}
 t_L=\Theta\!\left(\sqrt{\alpha+L^{-2}}\right).
 \label{eq:spider-prefix-t-asymptotic}
\end{equation}
Since
$\operatorname{vol}(U_L)=B(2L+1)$, a known-face semantic solve to tolerance
$\tau$ costs
\begin{equation}
 W_{B,L}=O\!\left(
  \frac{B(2L+1)}{\sqrt{\alpha+L^{-2}}}
  \log\frac{2}{\tau\sqrt{\alpha+L^{-2}}}
 \right).
 \label{eq:spider-prefix-sor-work}
\end{equation}
\end{corollary}

\begin{proof}
Substitute \eqref{eq:spider-prefix-cross-singular} into
\eqref{eq:fixed-face-optimal-sor}.  Moreover,
\begin{equation*}
 t_L^2=(1-c_\alpha^2)
 +c_\alpha^2\sin^2\!\left(\frac{\pi}{2(L+1)}\right).
\end{equation*}
Here $1-c_\alpha^2=4\alpha/(1+\alpha)^2$ and the sine term is
$\Theta(L^{-2})$, proving \eqref{eq:spider-prefix-t-asymptotic}.  Every
selected arm vertex has ambient degree two because $L<N$, while the hub has
degree $B$, giving $B+2BL=B(2L+1)$ and then
\eqref{eq:spider-prefix-sor-work}.
\end{proof}

This formula has a useful interpretation:
\begin{equation}
 \frac1{t_L}
 =\Theta\!\left(\min\{L,\alpha^{-1/2}\}\right).
 \label{eq:spider-two-regimes}
\end{equation}
Before the arm reaches the diffusion length $\alpha^{-1/2}$, its length is
the condition scale; afterward, teleportation caps the condition number and
the square-root PageRank scale takes over.  The spider does not create an
$\alpha^{-2}$ condition number.  Its difficulty is repeated volume and
support evolution, not worse fixed-face conditioning.

For the complete finite spider rather than a proper Dirichlet prefix,
$\sigma_U=1$ because every connected graph's normalized adjacency has the
degree eigenvector.  Thus
\begin{equation}
 t_{\rm full}=\frac{2\sqrt\alpha}{1+\alpha},
 \qquad
 W_{B,N}=O\!\left(
  \frac{BN}{\sqrt\alpha}
  \log\frac{2}{\tau\sqrt\alpha}
 \right).
 \label{eq:full-spider-sor-work}
\end{equation}

\subsection{Exact radial PPR profile}

For a hub seed, symmetry makes the degree-normalized PPR value
\begin{equation}
 y_j:=\frac{\pi_v}{d_v}
 \qquad (v\text{ at depth }j)
 \label{eq:spider-radial-y}
\end{equation}
constant on each level.  Put
\begin{equation}
 \lambda_\alpha:=\frac{1-\sqrt\alpha}{1+\sqrt\alpha}.
 \label{eq:spider-lambda}
\end{equation}
Solving the one-dimensional recurrence gives
\begin{align}
 y_0&=\frac{\sqrt\alpha}{B}
       \frac{1+\lambda_\alpha^{2N}}
            {1-\lambda_\alpha^{2N}},
 \label{eq:spider-center-profile}\\
 y_j&=y_0
 \frac{\lambda_\alpha^j+\lambda_\alpha^{2N-j}}
      {1+\lambda_\alpha^{2N}},
 \qquad 1\leq j\leq N.
 \label{eq:spider-arm-profile}
\end{align}
In particular, $0<y_j\leq y_0$.

\begin{proof}[Derivation]
The mass PageRank equation divided by degree is
\begin{equation*}
 y_v=\frac{2\alpha}{1+\alpha}\frac{s_v}{d_v}
 +c_\alpha\frac1{d_v}\sum_{u\sim v}y_u.
\end{equation*}
Away from the hub this is
$y_j=(c_\alpha/2)(y_{j-1}+y_{j+1})$, whose characteristic roots are
$\lambda_\alpha$ and $\lambda_\alpha^{-1}$.  The endpoint condition
$y_N=c_\alpha y_{N-1}$ selects a reflected profile proportional to
$\lambda_\alpha^j+\lambda_\alpha^{2N-j}$.  The hub condition
$y_0=2\alpha/((1+\alpha)B)+c_\alpha y_1$ yields
\eqref{eq:spider-center-profile}.  Finally,
$\lambda^j+\lambda^{2N-j}\leq1+\lambda^{2N}$ for $0<\lambda<1$.
\end{proof}

Equations \eqref{eq:spider-center-profile}--\eqref{eq:spider-arm-profile}
identify a natural output-scale family.  Take
\begin{equation}
 N=\Theta(\alpha^{-1/2}),
 \qquad
 B=\Theta(\sqrt\alpha/\varepsilon_{\rm ppr}).
 \label{eq:spider-output-scale-family}
\end{equation}
in the nontrivial regime
$\varepsilon_{\rm ppr}=O(\sqrt\alpha)$, with integer rounding.
Then $\max_j y_j=\Theta(\varepsilon_{\rm ppr})$ and
$BN=\Theta(1/\varepsilon_{\rm ppr})$.  A constant-factor reduction of the
zero-start semantic error in $O(1/\sqrt\alpha)$ sweeps would therefore give
the desired log-free work
\begin{equation}
 O\!\left(\frac1{\sqrt\alpha\,\varepsilon_{\rm ppr}}\right).
 \label{eq:spider-desired-log-free}
\end{equation}

The fixed-face two-norm theorem proves the same expression with a logarithm,
but it does not prove the constant-factor semantic statement because the
two-norm conversion can lose a factor depending on the number of arms.  The
following maximum-norm wave argument removes that loss.

\subsection{Exact log-free semantic damping for plain SOR}

Let $x^{(k)}$ be the zero-start iterate after $k$ complete red--black sweeps,
where every even-depth shell is relaxed first and every odd-depth shell
second.  On the complete spider use the graph-global parameter
\begin{equation}
 \omega_\alpha=1+\lambda_\alpha^2,
 \qquad
 \zeta_\alpha:=\lambda_\alpha^2.
 \label{eq:radial-sor-parameters}
\end{equation}

\begin{theorem}[Log-free semantic damping on every finite symmetric spider]
\label{thm:radial-semantic-damping}
For every $0<\alpha<1$, every number of arms $B\geq1$, every arm length
$N\geq1$, and every integer $k\geq0$,
\begin{align}
 \max_v\frac{|x_v^{(k)}-x_v^0|}{\sqrt{d_v}}
 &\leq
 \zeta_\alpha^k
 \left(1+2k(1-\zeta_\alpha)\right)
 \max_v\frac{x_v^0}{\sqrt{d_v}}
 \label{eq:radial-semantic-exact-envelope}\\
 &\leq
 8(1+k\sqrt\alpha)e^{-4k\sqrt\alpha}
 \max_v\frac{x_v^0}{\sqrt{d_v}}.
 \label{eq:radial-semantic-exponential-envelope}
\end{align}
The constants are independent of both $B$ and $N$.
\end{theorem}

\begin{proof}
Write
\begin{equation*}
 e_v^{(k)}:=
 \frac{x_v^0-x_v^{(k)}}{\sqrt{d_v}}.
\end{equation*}
Dividing the homogeneous coordinate-SOR error equation by $\sqrt{d_v}$
gives the semantic update
\begin{equation}
 e_v^+=(1-\omega_\alpha)e_v
 +\frac{\omega_\alpha c_\alpha}{d_v}
   \sum_{u\sim v}e_u
 =-\zeta_\alpha e_v
 +\frac{2\lambda_\alpha}{d_v}
   \sum_{u\sim v}e_u.
 \label{eq:radial-semantic-coordinate-update}
\end{equation}
Hub seeding and equal arm lengths preserve radial symmetry.  Denote the
common depth-$j$ value by $e_j^{(k)}$.

Reflect the depth sequence evenly at $0$ and $N$ and then extend it with
period $2N$.  This extension turns the endpoint averages into ordinary
nearest-neighbor averages on a cycle.  For indices modulo $N$, put
\begin{equation*}
 E_i^{(k)}:=e_{2i}^{(k)},
 \qquad
 O_i^{(k)}:=e_{2i+1}^{(k)},
\end{equation*}
and let $L$ be the cyclic shift on these $N$ cells.  Equation
\eqref{eq:radial-semantic-coordinate-update} becomes
\begin{align}
 E^{(k+1)}
 &=-\zeta_\alpha E^{(k)}
   +\lambda_\alpha(I+L^{-1})O^{(k)},
 \label{eq:radial-even-update}\\
 O^{(k+1)}
 &=-\zeta_\alpha O^{(k)}
   +\lambda_\alpha(I+L)E^{(k+1)}.
 \label{eq:radial-odd-update}
\end{align}
Eliminating the opposite color shows that either
$Z^{(k)}=E^{(k)}$ or $Z^{(k)}=O^{(k)}$ satisfies
\begin{equation}
 Z^{(k+1)}
 =2\zeta_\alpha MZ^{(k)}
  -\zeta_\alpha^2Z^{(k-1)},
 \qquad
 M:=\frac{L+L^{-1}}2.
 \label{eq:radial-chebyshev-recurrence}
\end{equation}
After dividing by $\zeta_\alpha^k$, the recurrence is the Chebyshev wave
equation.  For $k\geq1$ its exact solution is
\begin{equation}
 \frac{Z^{(k)}}{\zeta_\alpha^k}
 =T_k(M)Z^{(0)}
 +U_{k-1}(M)
   \left(\frac{Z^{(1)}}{\zeta_\alpha}-MZ^{(0)}\right).
 \label{eq:radial-chebyshev-solution}
\end{equation}
Because $M=(L+L^{-1})/2$,
\begin{align}
 T_k(M)&=\frac{L^k+L^{-k}}2,
 \label{eq:radial-chebyshev-first-kind}\\
 U_{k-1}(M)&=
 \sum_{q=0}^{k-1}L^{k-1-2q}.
 \label{eq:radial-chebyshev-second-kind}
\end{align}
Every shift is an isometry in maximum norm.  Therefore the two operator norms
in \eqref{eq:radial-chebyshev-first-kind}--%
\eqref{eq:radial-chebyshev-second-kind} are at most $1$ and $k$,
respectively.  This is the dimension-free step that an Euclidean spectral
bound misses.

It remains to bound the two initial wave mismatches.  At zero start,
$e_j^{(0)}=y_j$, so \eqref{eq:spider-arm-profile} gives
\begin{equation}
 \|E^{(0)}\|_\infty=y_0,
 \qquad
 \|O^{(0)}\|_\infty\leq y_0.
 \label{eq:radial-initial-color-norms}
\end{equation}
The source lies on the even color.  The odd PPR equations give
\begin{equation*}
 O^{(0)}=\frac{c_\alpha}{2}(I+L)E^{(0)}.
\end{equation*}
Using $c_\alpha/\lambda_\alpha=2/(1+\zeta_\alpha)$ in
\eqref{eq:radial-even-update} yields
\begin{equation}
 \frac{E^{(1)}}{\zeta_\alpha}-ME^{(0)}
 =s_\alpha(I+M)E^{(0)},
 \qquad
 s_\alpha:=\frac{1-\zeta_\alpha}{1+\zeta_\alpha}.
 \label{eq:radial-even-wave-mismatch}
\end{equation}
Hence its maximum norm is at most
$2(1-\zeta_\alpha)y_0$.

For the odd color, direct substitution of
\eqref{eq:radial-even-update} into
\eqref{eq:radial-odd-update} gives
\begin{equation}
 \frac{O^{(1)}}{\zeta_\alpha}-MO^{(0)}
 =(M-\zeta_\alpha I)O^{(0)}.
 \label{eq:radial-odd-wave-mismatch}
\end{equation}
The explicit profile is decreasing from either reflected endpoint and, for
$0\leq j\leq N-2$, satisfies
\begin{equation}
 0\leq y_j-y_{j+2}
 =\frac{(1-\zeta_\alpha)y_0}
        {1+\lambda_\alpha^{2N}}
   \left(
    \lambda_\alpha^j
    -\lambda_\alpha^{2N-j-2}
   \right)
 \leq(1-\zeta_\alpha)y_0.
 \label{eq:radial-profile-two-step-difference}
\end{equation}
The reflected boundary differences obey the same bound.  Consequently,
$\|(M-I)O^{(0)}\|_\infty\leq
(1-\zeta_\alpha)y_0$, and
\eqref{eq:radial-odd-wave-mismatch} also has norm at most
$2(1-\zeta_\alpha)y_0$.

Applying \eqref{eq:radial-chebyshev-solution} to both colors proves
\eqref{eq:radial-semantic-exact-envelope}.  Finally,
\begin{equation*}
 1-\zeta_\alpha
 =\frac{4\sqrt\alpha}{(1+\sqrt\alpha)^2}
 \leq4\sqrt\alpha,
 \qquad
 \zeta_\alpha^k=\lambda_\alpha^{2k}
 \leq e^{-4k\sqrt\alpha}.
\end{equation*}
Since $1+8z\leq8(1+z)$ for $z\geq0$, the exponential envelope follows.
\end{proof}

\begin{proposition}[The source-first sweep order is necessary]
\label{prop:radial-sweep-order-witness}
The sharp envelope
\eqref{eq:radial-semantic-exact-envelope} does not hold unchanged when the
source is on the second color and the original even--odd sweep order is kept.
Already on the two-vertex path, take $\alpha=1/9$, start from zero, put the
source on the odd vertex, and relax the even vertex first.  After one full
sweep the maximum semantic error has fallen by only $4/5$, whereas the
right-hand-side factor in
\eqref{eq:radial-semantic-exact-envelope} is $5/8$.
\end{proposition}

\begin{proof}
Here $\lambda_\alpha=1/2$ and $\zeta_\alpha=1/4$.  The exact semantic PPR
error at zero start is $(4/9,5/9)$, ordered as even then odd.  The two
coordinate updates from
\eqref{eq:radial-semantic-coordinate-update} give
\begin{equation*}
 e_{\rm even}^{(1)}
 =-\frac14\frac49+\frac59=\frac49,
 \qquad
 e_{\rm odd}^{(1)}
 =-\frac14\frac59+\frac49=\frac{11}{36}.
\end{equation*}
Thus the maximum-error ratio is $(4/9)/(5/9)=4/5$.  In contrast, the claimed
one-sweep factor would be
$\zeta_\alpha(1+2(1-\zeta_\alpha))=5/8$.  Swapping the two colors puts the
source first and restores the theorem by symmetry.
\end{proof}

\begin{corollary}[Log-free output-scale plain SOR]
\label{cor:radial-output-scale-sor}
On the family \eqref{eq:spider-output-scale-family}, choose the hidden
constants so that the zero-start semantic error is a fixed constant multiple
of $\varepsilon_{\rm ppr}$.  A predetermined
$k=O(1/\sqrt\alpha)$ full-sweep count reaches semantic error
$\varepsilon_{\rm ppr}$ in
\begin{equation}
 W_{\rm radial\text{-}SOR}
 =O\!\left(
  \frac1{\sqrt\alpha\,\varepsilon_{\rm ppr}}
 \right)
 \label{eq:radial-output-scale-sor-work}
\end{equation}
charged degree work, including every sweep and explicit output.
\end{corollary}

\begin{proof}
Theorem~\ref{thm:radial-semantic-damping} reduces any fixed initial-to-target
ratio in $O(1/\sqrt\alpha)$ sweeps.  A full spider sweep costs its degree
volume $2BN=\Theta(1/\varepsilon_{\rm ppr})$, and the output has
$1+BN=O(1/\varepsilon_{\rm ppr})$ coordinates.
\end{proof}

This closes the former plain-SOR radial conjecture, but only for a complete
equal-arm spider, a hub seed, zero start, full red--black sweeps, and the
predetermined semantic schedule, with the hub color relaxed first.
Proposition~\ref{prop:radial-sweep-order-witness} shows that the order clause
cannot simply be deleted.  Unequal arms, nonradial seeds, support discovery,
and changing faces are not consequences of the theorem.

\subsection{How the radial two-rung theorem differs}

The proof-owning \texttt{two\_rung\_direct\_theory} note gives a different
accelerated radial algorithm.  On a symmetric spider with $R$ arms, it groups
every depth shell into one commuting block, uses optimal SOR while the
normalized residual is above a coarse threshold, and finishes with unit
relaxation.  If
$B_{\rm edge}>1$ denotes the source theorem's admissibility constant, the
imported theorem assumes
\begin{equation*}
 1<B_{\rm rung}<B_{\rm edge},
 \qquad
 B_{\rm rung}R\varepsilon_{\rm ppr}<1,
\end{equation*}
and gives
\begin{equation}
 W_{\rm radial\text{-}two\text{-}rung}
 =O\!\left(
  \mathcal V_{\rm exp}
  \left[1+
   \frac{\log(1/(B_{\rm rung}R\varepsilon_{\rm ppr}))}
        {\sqrt\alpha}
  \right]
 \right)
 \label{eq:imported-radial-two-rung}
\end{equation}
when the arms are not yet exhausted.  Under the same two parameter
hypotheses, a second imported theorem handles finite arms: if the wave reaches
the far leaves, it certifies the closed pendant paths, Schur-deflates them,
solves the remaining center scalar, and reverse-recovers all arms within the
same asymptotic bound.

This imported result and
\Cref{thm:radial-semantic-damping} should not be merged.  The imported method
uses the note-scoped residual gate
$\tau=\alpha\varepsilon_{\rm ppr}$, a two-rung schedule, refreshed radial
selection, and possibly exact Schur deflation; it can explore only the needed
prefix.  The new theorem instead scans the complete spider every sweep and
uses a predetermined semantic count.  Both are accelerated, but their
locality and stopping contracts differ.

\subsection{A stronger end-to-end spider algorithm already exists}

For RPPR, forcing SOR to perform support discovery is unnecessary on a
spider.  The proof-owning \texttt{propagate\_settle\_framework} note gives an
exact kinetic response algorithm for a rooted spider with arbitrary arm
lengths.  Eliminating each active arm prefix leaves an arrowhead system.  Its
hub value is one monotone scalar, and every next arm event has a stable
threshold.  A heap reports each event once and reverse substitution recovers
the exact solution.

The imported charged bound is
\begin{equation}
 W_{\rm kinetic}
 =O\!\left(\operatorname{cvol}(S^\rho)
            \log(2+B)\right),
 \label{eq:kinetic-spider-import}
\end{equation}
where $S^\rho=\operatorname{supp}(x^\rho)$ and
$\operatorname{cvol}(S)=\sum_{v\in S}(1+d_v)$.  Since
$\operatorname{vol}(S^\rho)\leq1/\rho$ and
$\operatorname{cvol}(S^\rho)\leq2/\rho$ for a nonempty support, choosing
$\rho=\varepsilon_{\rm ppr}$ and applying
\Cref{thm:rppr-ppr-bias} gives
\begin{equation}
 W_{\rm kinetic}
 =O\!\left(
  \frac{\log(2+B)}{\varepsilon_{\rm ppr}}
 \right)
 \label{eq:kinetic-spider-ppr-work}
\end{equation}
for semantic PPR output in exact arithmetic.
If $S^\rho$ is empty, the source algorithm instead returns the zero vector
after its constant-work initial test.

This is stronger than \eqref{eq:spider-desired-log-free} in its
$\alpha$-dependence.  It is not evidence that the SOR analysis is wrong; it
uses a stronger persistent-response primitive that exploits the spider's
one-articulation geometry.  The algorithmic lesson is to use the response
solver on spiders and use SOR as a benchmark or a numerical tail, rather than
insist that one primitive be universal.
